\documentclass[10pt,a4paper]{amsart}
\usepackage[margin=19mm]{geometry}
\usepackage{amsmath,amssymb,mathtools}
\usepackage[scr=rsfs,cal=euler]{mathalfa}
\usepackage{xfrac}
\usepackage[unicode,hidelinks,bookmarks=false]{hyperref}
\newcommand{\ie}{i.e.\ }
\newcommand{\cf}{cf.\ }
\newcommand{\resp}{respectively\ }
\newcommand{\Subsection}{\subsection}
\providecommand{\nobreakdash}{\nobreak\hbox{-}}
\setlength{\emergencystretch}{2em}
\multlinegap0pt
\usepackage{amssymb}
\usepackage[shortlabels]{enumitem}
\usepackage{tensor}
\usepackage{braket}
\newtheorem{lemma}{Lemma}
\newcommand{\Cc}{\mathbf{C}}
\newcommand{\Qq}{\mathbf{Q}}
\newcommand{\fleche}[1]{\stackrel{#1}{\lra}}
\newcommand{\lra}{\longrightarrow}
\newcommand{\mfm}{\mathfrak{m}}
\def\setminus{\mathchoice
    {\mathbin{\vrule height .62ex width 1.61ex depth -.38ex}}% 12
    {\mathbin{\vrule height .62ex width 1.61ex depth -.38ex}}% 12
    {\mathbin{\vrule height .50ex width 0.85ex depth -.28ex}}%  9
    {\mathbin{\vrule height .20ex width 0.570ex depth -.24ex}}%  7
}
\title{Jacobian graphs}
\author{Arthur Forey, Javier Fres\'an, Emmanuel Kowalski, Yuval Wigderson}
\date{Source: arXiv:2603.13198v1 (CC BY 4.0)}
\begin{document}
\maketitle

\noindent\emph{Source and licence.} Jacobian graphs. Version arXiv:2603.13198v1 (CC BY 4.0), distributed by arXiv under the Creative Commons Attribution 4.0 International licence, http://creativecommons.org/licenses/by/4.0/, as recorded in the arXiv licence metadata for this exact version and shown on the abstract page https://arxiv.org/abs/2603.13198v1. Full licence notice: \texttt{licenses/arxiv-2603.13198-CC-BY-4.0.txt}.\par\smallskip

\noindent\textit{Editorial scope.} Counted unit: the article's lemma lm-nonzero (source0.tex lines 2058--2065). The article's complete original proof of it is delivered with it (source0.tex lines 2067--2124, one \texttt{proof} environment of 58 lines). No mathematics is written, repaired or re-derived: the statement and the proof are the article's own lines, in the article's own notation, and no retained line is substituted or re-flowed.  No further source-local context is needed: every cross-reference of every retained line resolves inside the two delivered spans.  Nothing was omitted from any retained span, so the package carries no omission record.  No retained line contains a citation, so the wrapper carries no bibliography and no bibliography entry.  No retained line contains a figure environment, an image-inclusion command or drawing code, so no image is delivered and the package declares no asset dependency.  Editorial formatting only, in addition: an AMS \texttt{amsart} preamble that restates the article's own declarations and macros used by the retained lines, namely the environments \texttt{lemma} on separate counters, together with the standard packages those lines need.  The retained environments print in this document as Lemma 1 in order of appearance, whereas the article prints the same items with the numbering of its own full document.  The complete native source of the article as distributed in the arXiv e-print for this exact version is bundled unchanged as \texttt{sources/196327/source0.tex}.
\begin{lemma}\label{lm-nonzero}
  Assume that $C$ has genus~$1$ and $\mfm$ is a double point. If $|k|$
  is coprime to $|C(k)|$, then the sum
  \[
    T(\chi)=\sum_{x\in (C\setminus \mfm)(k)}\chi(i_{\delta}(x))
  \]
  is non-zero for all characters~$\chi$ of $J_{\mfm}(k)$.
\end{lemma}

\begin{proof}
  We denote by~$p$ the characteristic of~$k$. We have the exact sequence
  \[
    1\to k\to J_{\mfm}(k)\fleche{\pi} C(k)\to 1.
  \]

  As before, we recall that the size of $C(k)$ is given by the formula
  \[
    |C(k)|=|k|+1-a_C(k),
  \]
  where the integer $a_C(k)$ is of the form
  \[
    a_C(k)=|k|^{1/2}(\alpha+\alpha^{-1})
  \]
  for some algebraic number~$\alpha$ with $|\alpha|=1$.
  
  Since we assume that the order of $C(k)$ is coprime to~$|k|$, the
  above exact sequence splits, and there exists a (section) homomorphism
  $\sigma\colon C(k)\to J_{\mfm}(k)$, defining an isomorphism
  \[
    J_{\mfm}(k)\to k\times C(k)
  \]
  by $x\mapsto (x-\sigma(\pi(x)),\pi(x))$.

  Let~$\chi$ be a character of $J_{\mfm}(k)$, and denote by $\chi_1$
  and~$\chi_2$ the characters of $k$ and $C(k)$ corresponding to $\chi$
  through this isomorphism; we also view these as characters of
  $J_{\mfm}(k)$ (e.g., we write $\chi_2(x)$ for $\chi_2(\pi(x))$ when
  $x\in J_{\mfm}(k)$).

  The values of~$\chi_1$ are $p$-th roots of unity. We write
  \[
    T(\chi)=\sum_{\xi^p=1}\xi\sum_{\substack{x\in (C\setminus
        \mfm)(k)\\\chi_1(i_{\delta}(x))=\xi}}\chi_2(i_{\delta}(x)),
  \]
  where~$\xi$ runs over $p$-th roots of unity in~$\Cc$. For each~$\xi$,
  the inner sum is a cyclotomic integer and belongs to the field
  generated by values of~$\chi_2$, which are $|C(k)|$-th roots of unity.
  Since~$p\nmid |C(k)|$, this field is linearly disjoint from
  $\Qq(e^{2i\pi/p})$, and it follows that $T(\chi)=0$ if and only if the
  cyclotomic integer
  \[
    \sum_{\substack{x\in (C\setminus
        \mfm)(k)\\\chi_1(i_{\delta}(x))=\xi}}\chi_2(i_{\delta}(x)),
  \]
  is \emph{independent} of~$\xi$. If we call this cyclotomic integer
  $\beta$, then we deduce that
  \[
    p\beta=\sum_{\xi^p=1} \sum_{\substack{x\in (C\setminus
        \mfm)(k)\\\chi_1(i_{\delta}(x))=\xi}}\chi_2(i_{\delta}(x))=
    \sum_{x\in (C\setminus
      \mfm)(k)}\chi_2(i_{\delta}(x))=-\chi_2(i_{\delta}(x_0)),
  \]
  where $\mfm=2(x_0)$ (since $i_{\delta}$ is a set-theoretic section
  of~$\pi$ on $C\setminus \mfm$). But this is absurd since the
  right-hand side is a unit whereas the left-hand side is not, and
  therefore $T(\chi)\not=0$.
\end{proof}

\end{document}
